\documentclass[12pt]{article}

\usepackage[utf8]{inputenc}
\usepackage{amsmath, amssymb, amsthm}
\usepackage{geometry}
\geometry{a4paper, margin=2.5cm}

\title{Formula lui Bayes: Definiție și Teoremă}
\author{}
\date{}

\theoremstyle{plain}
\newtheorem{theorem}{Teoremă}

\begin{document}

\maketitle

\section{Definiție}

Fie $(\Omega, \mathcal{F}, \mathbb{P})$ un spațiu de probabilitate.
Pentru două evenimente $A$ și $B$ cu $\mathbb{P}(B) > 0$, 
\textbf{probabilitatea condiționată} a lui $A$ dat $B$ este definită prin:



\[
\mathbb{P}(A \mid B)
= \frac{\mathbb{P}(A \cap B)}{\mathbb{P}(B)}.
\]



Această definiție stă la baza formulei lui Bayes.

\section{Formula lui Bayes}

\begin{theorem}[Formula lui Bayes]
Fie $A$ și $B$ două evenimente cu $\mathbb{P}(A) > 0$ și $\mathbb{P}(B) > 0$.
Atunci probabilitatea condiționată a lui $A$ dat $B$ este:



\[
\mathbb{P}(A \mid B)
= \frac{\mathbb{P}(B \mid A)\,\mathbb{P}(A)}
       {\mathbb{P}(B)}.
\]



Dacă evenimentul $B$ poate fi descompus printr-o partiție finită
$\{A_1, A_2, \dots, A_n\}$ a lui $\Omega$, atunci:



\[
\mathbb{P}(A_k \mid B)
= \frac{\mathbb{P}(B \mid A_k)\,\mathbb{P}(A_k)}
       {\sum_{i=1}^n \mathbb{P}(B \mid A_i)\,\mathbb{P}(A_i)}.
\]



Aceasta este forma generală a teoremei lui Bayes.
\end{theorem}

\section{Interpretare}

Formula lui Bayes permite actualizarea probabilităților în funcție de informații noi.
Probabilitatea \emph{a priori} $\mathbb{P}(A)$ este actualizată în probabilitatea
\emph{a posteriori} $\mathbb{P}(A \mid B)$ după observarea evenimentului $B$.

\end{document}
